\documentclass[
  10pt,
  ngerman,
  a4paper,
  toc=graduated,
  numbers=noenddot,
  headings=standardclasses,
  nonumberlist
]{scrartcl}

\setlength{\parindent}{0pt}
\setlength{\parskip}{0.35\baselineskip plus 0.5pt minus 0.5pt}
\setlength{\emergencystretch}{2em}

\usepackage[utf8]{inputenc}
\usepackage[T1]{fontenc}
\usepackage{tgheros}
\renewcommand{\familydefault}{\sfdefault}
\usepackage{babel}
%\usepackage{lmodern}
\usepackage[top=2cm, right=2.5cm, bottom=2cm, left=2.5cm, footskip=0.5cm]{geometry}
\usepackage{calc}
\usepackage{graphicx}
\usepackage{svg}
\usepackage{amsmath}
\usepackage{tabularx}
\usepackage{booktabs}
\usepackage[official]{eurosym}
\usepackage{lastpage}
\usepackage{etoolbox}
\usepackage{amssymb}
\usepackage{enumitem}
\usepackage{float}
%\usepackage{natbib}
%\bibliographystyle{agsm}
%\bibliographystyle{IEEETran}
\newcommand{\code}[1]{\nolinkurl{#1}}
\usepackage[hidelinks]{hyperref}
\usepackage{dramatist}
\usepackage{colortbl}
\usepackage{wrapfig}

\usepackage[table]{xcolor}
\usepackage{multirow}
\usepackage{array}
\usepackage{comment}

\definecolor{ihk_blue}{RGB}{28, 54, 97}
\definecolor{hkblue}{RGB}{32,58,102}      % dunkles HK-Blau
\definecolor{hklight}{RGB}{220,226,234}   % helles Blau-Grau
\definecolor{hkline}{RGB}{55,65,80}       % dunkle Tabellenlinien

% Prozentangaben, z. B. \pct{62{,}5} -> 62,5 %
\newcommand{\pct}[1]{#1\,\%}

% Alias für die in den Tabellen verwendete Farbe
\colorlet{lightblue}{hklight}
\colorlet{tablehead}{hkblue}
\colorlet{tablelight}{hklight}
\renewcommand{\arraystretch}{1.0}
\setlength{\arrayrulewidth}{0.65pt}
\arrayrulecolor{hkline}

\usepackage{adjustbox,lipsum}
\usepackage{makecell}
\usepackage{caption}
\usepackage{longtable}
\usepackage{tabularx}
\usepackage{ragged2e}
\usepackage{array}
\newcolumntype{Y}{>{\RaggedRight\arraybackslash}X}
\newcolumntype{L}[1]{>{\RaggedRight\arraybackslash}p{#1}}
\newcolumntype{R}[1]{>{\RaggedLeft\arraybackslash}p{#1}}
\newcolumntype{C}[1]{>{\Centering\arraybackslash}p{#1}}

% ------------------------------------------------------------
% Code-Listings
% ------------------------------------------------------------

% Nur einmal laden
\usepackage{listings}
\usepackage[scaled=0.95]{inconsolata}

% Codefarben
\definecolor{codebg}{RGB}{245,247,250}
\definecolor{codeframe}{RGB}{210,214,220}
\definecolor{codenumber}{RGB}{130,130,130}
\definecolor{codekeyword}{RGB}{28,54,97}
\definecolor{codestring}{RGB}{20,120,60}
\definecolor{codecomment}{RGB}{110,120,130}
\definecolor{codeclass}{RGB}{125,55,160}
\definecolor{codefunction}{RGB}{35,95,160}

% Einheitlicher Grundstil
\lstdefinestyle{basecode}{
    basicstyle=\ttfamily\fontsize{10pt}{11.5pt}\selectfont,
    numbers=none,
    numberstyle=\tiny\color{codenumber},
    stepnumber=1,
    numbersep=7pt,
    frame=single,
    framerule=0.5pt,
    rulecolor=\color{codeframe},
    backgroundcolor=\color{codebg},
    showstringspaces=false,
    showtabs=false,
    showspaces=false,
    breaklines=true,
    breakatwhitespace=true,
    keepspaces=true,
    columns=fullflexible,
    tabsize=4,
    captionpos=t,
    aboveskip=1em,
    belowskip=1em,
    resetmargins=true,
    xleftmargin=0pt,
    xrightmargin=0pt,
    literate=
        {ä}{{\"a}}1
        {ö}{{\"o}}1
        {ü}{{\"u}}1
        {Ä}{{\"A}}1
        {Ö}{{\"O}}1
        {Ü}{{\"U}}1
        {ß}{{\ss}}1,
    framesep=5pt
}
% Python
\lstdefinelanguage{PythonCustom}{
    keywords={
        False,None,True,and,as,assert,async,await,break,class,
        continue,def,del,elif,else,except,finally,for,from,global,
        if,import,in,is,lambda,nonlocal,not,or,pass,raise,return,
        try,while,with,yield
    },
    sensitive=true,
    comment=[l]{\#},
    morestring=[b]',
    morestring=[b]",
    morestring=[s]{'''}{'''},
    morestring=[s]{"""}{"""},
    keywordstyle=\bfseries\color{codekeyword},
    commentstyle=\itshape\color{codecomment},
    stringstyle=\color{codestring}
}

\lstdefinestyle{pythoncode}{
    style=basecode,
    language=PythonCustom,
    emph={
        Agent,App,DocumentServiceClient,SearchServiceClient,
        ImportDocumentsRequest,GcsSource,ReconciliationMode,
        ParsedQuestion,Client,Part
    },
    emphstyle=\bfseries\color{codeclass},
    emph={[2]
        get_questions,starte_interview,
        speichere_antwort_und_naechste_frage,
        zeige_interview_status,generiere_ergebnisdateien,
        build_result_documents,build_result_docx,
        build_transcript_docx,save_artifact,
        import_documents,result,append,sort,strip,get
    },  
    emphstyle={[2]\color{codefunction}}
}

% JSON
\lstdefinelanguage{JsonCustom}{
    morestring=[b]",
    stringstyle=\color{codestring},
    showstringspaces=false,
    literate=
     *{0}{{{\color{black}0}}}{1}
      {1}{{{\color{black}1}}}{1}
      {2}{{{\color{black}2}}}{1}
      {3}{{{\color{black}3}}}{1}
      {4}{{{\color{black}4}}}{1}
      {5}{{{\color{black}5}}}{1}
      {6}{{{\color{black}6}}}{1}
      {7}{{{\color{black}7}}}{1}
      {8}{{{\color{black}8}}}{1}
      {9}{{{\color{black}9}}}{1}
      {:}{{{\color{codekeyword}{:}}}}{1}
      {,}{{{\color{codekeyword}{,}}}}{1}
      {\{}{{{\color{codekeyword}{\{}}}}{1}
      {\}}{{{\color{codekeyword}{\}}}}}{1}
      {[}{{{\color{codekeyword}{[}}}}{1}
      {]}{{{\color{codekeyword}{]}}}}{1}
}

\lstdefinestyle{jsoncode}{
    style=basecode,
    language=JsonCustom
}

% JavaScript
\lstdefinelanguage{JavaScriptCustom}{
    keywords={
        function,return,const,let,var,if,else,for,while,do,
        switch,case,break,continue,new,true,false,null,undefined,
        async,await,try,catch,finally,throw,class,extends,import,
        export,from,default,this,of,in
    },
    sensitive=true,
    comment=[l]{//},
    morecomment=[s]{/*}{*/},
    morestring=[b]",
    morestring=[b]',
    keywordstyle=\bfseries\color{codekeyword},
    commentstyle=\itshape\color{codecomment},
    stringstyle=\color{codestring}
}

\lstdefinestyle{jscode}{
    style=basecode,
    language=JavaScriptCustom,
    emph={
        normalizeText,tokenize,levenshtein,tokenJaccard,
        scoreCandidate,resolveKey,searchSubstitutes
    },
    emphstyle=\color{codefunction}
}


% ------------------------------------------------------------
% Ende Code-Listings
% ------------------------------------------------------------
\usepackage[acronym, toc]{glossaries}

\newacronym{erd}{ERD}{Entity-Relation-Diagram}
\newacronym{mvp}{MVP}{Minimum Viable Product}
\newacronym{jav}{JAV}{Jugend- und Auszubildendenvertretung}
\newacronym{br}{BR}{Betriebsrat}
\newacronym{dod}{DoD}{Definition of Done}
\newacronym{nps}{NPS}{Net Promoter Score}
\newacronym{gcs}{GCS}{Google Cloud Storage}
\newacronym{adk}{ADK}{Agent Development Kit}
\newacronym{ki}{KI}{Künstliche Intelligenz}
\newacronym{ui}{UI}{User Interface}

%\newglossaryentry{gitlabrunner}{name={GitLab-Runner},description={Anwendung, die auf Aktionen eines Git Repositories reagiert und vorkonfigurierte Jobs ausführt}}
%\newglossaryentry{rest}{name={REST}, description = {Representational State Transfer, Ein Paradigma für die Architektur von \gls{crud} Webanwendungen }}
\newglossarystyle{abkuerzungen}{
  \renewenvironment{theglossary}
    {\begin{longtable}{@{}p{3.0cm}@{}p{\dimexpr\textwidth-3.0cm\relax}@{}}}
    {\end{longtable}}

  \renewcommand*{\glossentry}[2]{%
    \glstarget{##1}{%
      \makebox[3.0cm][l]{\textbf{\glossentryname{##1}}\dotfill}%
    }&
    \glossentrydesc{##1} \\
  }

  \renewcommand*{\subglossentry}[3]{}
  \renewcommand*{\glsgroupskip}{}
}

% ------------- Abbildungen mit Tikz erstellen
\usepackage{tikz}
\usetikzlibrary{arrows,automata,shapes,calc}


% ------------- Zeilenabstand
\usepackage{setspace}
\singlespacing

% ------------- Kopfzeile
\usepackage[automark]{scrlayer-scrpage}
\pagestyle{scrheadings}
\clearpairofpagestyles
\ihead{}
%\ihead[\rightmark]{\pagemark}
\chead{}
%\chead{}{}
%\ohead{\pagemark}
%\ohead[\pagemark]{\rightmark}
\ofoot{\pagemark}
%\ofoot[\pagemark]{\rightmark}


\renewcommand{\contentsname}{Inhaltsverzeichnis}
\renewcommand{\listfigurename}{Abbildungsverzeichnis}
\renewcommand{\listtablename}{Tabellenverzeichnis}

\RedeclareSectionCommand[
  font=\normalfont\bfseries
]{section}

\RedeclareSectionCommand[
  font=\normalfont\bfseries
]{subsection}

\RedeclareSectionCommand[
  font=\normalfont\bfseries
]{subsubsection}

%-------------- Ende allgemeine Angaben -----------------------------
\begin{document}

%-------------- Beginn Titelseite -----------------------------------
\hypersetup{pageanchor=false}
\thispagestyle{empty}

\vspace{1cm}

\begin{figure}[htbp]
    \centering

    \begin{minipage}[t]{0.45\textwidth}
        \raggedright
        \includesvg[width=5cm,height=10cm,keepaspectratio]{images/bhh_logo}
    \end{minipage}
    \hfill
    \begin{minipage}[t]{0.45\textwidth}
        \raggedleft
        \includesvg[width=5cm,height=10cm,keepaspectratio]{images/signal-iduna-logo}
    \end{minipage}

\end{figure}

\vspace{1cm}
\begin{center}
	{\Large Praxisvalidierungsarbeit III} \\
	\vspace{3cm}
	{\huge \textbf{Entwicklung eines KI-Agenten zur Sicherung impliziten Wissens}}\\
    \vspace{0.3cm}
    %{\LARGE \textbf{{Umsetzung mit Google \glslink{adk}{ADK} und Vertex AI}}}\\
    {\LARGE \textbf{{Empirische Einzelfallanalyse und synthetische Validierung des Bewertungsverfahrens}}}\\
    
\end{center}

\vspace{1.5cm}
\begin{center}
	Abgabe: Hamburg, den 11.08.2026\\[3\baselineskip]
	\textbf{\textcolor{ihk_blue}{Matrikelnummer: 236913}}

    Bachelor of Science (B.Sc.) \\[4\baselineskip]
	
	Erstbetreuung \\
    Prof. Dr. Henning Klaffke \\[2\baselineskip]

    Zweitbetreuung \\
    Prof. Dr. Robert Mertens \\

\end{center}
%-------------- Ende Titelseite -------------------------------------


% ------------------------------------------------------------
% Verzeichnisse und Erklärungen
% ------------------------------------------------------------
\newpage
\hypersetup{pageanchor=true}
\pagenumbering{Roman}
\tableofcontents

\newpage
\section*{Sperrvermerk}
\addcontentsline{toc}{section}{Sperrvermerk}
Die vorliegende Praxisvalidierungsarbeit enthält vertrauliche interne Informationen der SIGNAL IDUNA Gruppe. Sie ist ausschließlich zur Bewertung im Rahmen des Studiums an der Beruflichen Hochschule Hamburg bestimmt. Eine Veröffentlichung, Vervielfältigung oder Weitergabe der Arbeit oder einzelner Inhalte an Dritte ist ohne ausdrückliche Zustimmung der SIGNAL IDUNA Gruppe nicht gestattet.

\vspace{2cm}
\begin{tabular}{@{}p{6cm}p{6cm}@{}}
\hrulefill & \hrulefill\\
Hamburg, den 11.08.2026 & Matrikelnummer 236913
\end{tabular}

\vspace{2cm}
\hrulefill\\
Unterschrift

\newpage
\section*{Selbstständigkeitserklärung}
\addcontentsline{toc}{section}{Selbstständigkeitserklärung}
Hiermit erkläre ich an Eides statt, dass ich die vorliegende wissenschaftliche Arbeit selbstständig verfasst und keine anderen als die angegebenen Hilfsmittel benutzt habe. Quellen, die im Wortlaut oder dem Sinn nach verwendet wurden, sind kenntlich gemacht.

Werkzeuge auf Basis künstlicher Intelligenz wurden unterstützend zur Strukturierung, sprachlichen Überarbeitung und Formulierung eingesetzt. Die inhaltliche Verantwortung, fachliche Prüfung und finale Bearbeitung der Inhalte liegen vollständig bei mir. In dieser Arbeit verwendete synthetische Testdaten und modellierte Annahmen werden ausdrücklich als solche gekennzeichnet.

\vspace{2cm}
\begin{tabular}{@{}p{6cm}p{6cm}@{}}
\hrulefill & \hrulefill\\
Hamburg, den 11.08.2026 & Matrikelnummer 236913
\end{tabular}

\vspace{2cm}
\hrulefill\\
Unterschrift

\newpage
\listoffigures
\addcontentsline{toc}{section}{Abbildungsverzeichnis}

\newpage
\listoftables
\addcontentsline{toc}{section}{Tabellenverzeichnis}

% ------------------------------------------------------------
% Hauptteil
% ------------------------------------------------------------
\newpage
\pagenumbering{arabic}
\setcounter{page}{1}

\section{Einleitung und Forschungsfrage}

Beim Ausscheiden, Rollenwechsel oder längeren Ausfall erfahrener Beschäftigter besteht das Risiko, dass nicht nur formale Zuständigkeiten, sondern auch informelle Abläufe, persönliche Netzwerke, Fehlererfahrungen und bewährte Vorgehensweisen verloren gehen. Gelashvili-Luik et al. zeigen in einer systematischen Literaturübersicht, dass KI-gestützte Wissensmanagementsysteme zwar neue Möglichkeiten zur Erfassung und Bereitstellung von Wissen eröffnen, ihr Erfolg aber ebenso von Datenqualität, organisatorischer Einbettung, Qualifikation, Vertrauen und menschlicher Kontrolle abhängt \cite{gelashvili2025}. Junghans et al. beschreiben einen thematisch nahen Ansatz, bei dem informelle betriebliche Dialoge mit Sprachverarbeitung und Sprachmodellen in strukturiertes Organisationswissen überführt werden \cite{junghans2026}.

\subsection{Gegenstand des zugrunde liegenden IHK-Praxisprojekts}

Die vorliegende Praxisvalidierungsarbeit baut auf dem IHK-Praxisprojekt \emph{Entwicklung eines KI-Agenten zur Sicherung impliziten Wissens -- Umsetzung mit Google ADK und Vertex AI} auf. Ziel dieses Projekts war die Entwicklung eines lauffähigen Minimum Viable Product (MVP), das einen bislang überwiegend manuellen und uneinheitlichen Wissenssicherungsprozess standardisiert. Der Agent richtet sich an Beschäftigte, deren fachliches, technisches oder prozessbezogenes Erfahrungswissen für spätere Übergaben, Einarbeitungen oder Vertretungssituationen gesichert werden soll \cite[S.~1--2]{sommer2026projekt}.

Der Agent führt ein textbasiertes, leitfadengestütztes Interview in zwei Abschnitten. Teil~1 enthält allgemeine und rollenbezogene Fragen zu Netzwerken, Prozessen, Werkzeugen und Ressourcen. Teil~2 enthält 13 projektspezifische Fragen und wird für jedes von der interviewten Person benannte kritische Projekt beziehungsweise Aufgabenfeld wiederholt. Die Fragen adressieren unter anderem wichtige Ansprechpartner, Kommunikationsbeziehungen, Hintergründe von Entscheidungen, bewältigte Fehler und Herausforderungen, typische Probleme, informelle Abhängigkeiten, Fehlerrisiken, Regeltätigkeiten, praktische Tipps, zukünftige Aufgaben und Trends, Dokumentationsorte sowie externe Informationsquellen. Jede Frage besitzt bereits im Fragenkatalog eine Ergebniskategorie, ein Keyword und ein definiertes Erkenntnisinteresse \cite[S.~2]{sommer2026projekt}.

Technisch wird der MVP als Single-Agent-Anwendung mit dem Google Agent Development Kit umgesetzt. Die nutzende Person interagiert über die ADK Web UI; Gemini~2.5~Pro wird über Vertex~AI für Dialog und sprachliche Verdichtung genutzt. Die vorbereiteten Fragen und ergänzende Kontextdokumente werden über die Discovery Engine bereitgestellt. Während des Interviews speichert die Anwendung Frage, Rohantwort, Kurzfassung, Kategorie und Keyword im Session State. Nach Abschluss werden ein Transkript, ein Ergebnisdokument für allgemeines rollenbezogenes Wissen und für jedes angegebene Thema ein separates projektspezifisches DOCX-Dokument erzeugt. Die Kurzfassungen werden auf zwei bis vier Sätze begrenzt und sollen keine neuen Informationen, Kategorien oder Keywords erzeugen \cite[S.~9--12]{sommer2026projekt}.

Für die Interpretation der PVA ist die Abgrenzung des MVP wesentlich. Nicht Bestandteil des IHK-Projekts waren eine Produktivsetzung, die automatische Ablage der Ergebnisdateien, eine Sprachschnittstelle, eine vollständige inhaltliche Konsistenzprüfung der Nutzendeneingaben und dynamische Rückfragen zur Klärung widersprüchlicher Angaben \cite[S.~2,~5]{sommer2026projekt}. Die PVA bewertet daher kein autonomes Wissensmanagementsystem, sondern ein prototypisches Artefakt, das den Interviewablauf und die strukturierte Aufbereitung unterstützt.

Die IHK-Projektdokumentation wird in dieser Arbeit als unveröffentlichte Primärquelle für Funktionsumfang, Testartefakte und Prozessannahmen zitiert. Sie wird aufgrund ihres Umfangs und des Sperrvermerks nicht vollständig im Anhang reproduziert. Stattdessen enthält der Anhang der PVA die für die Evaluation unmittelbar benötigten Daten und Ableitungen, insbesondere Referenzstandard, Codierung und synthetische Testdefinitionen.

\subsection{Ziel der Praxisvalidierungsarbeit und Forschungsfrage}

Die ursprüngliche Test- und Abnahmephase überprüfte vor allem die funktionalen Must-have-Anforderungen des MVP. Insgesamt wurden vier lokale Agentendurchläufe mit vergleichbaren Rahmenbedingungen durchgeführt; im IHK-Bericht sind zwei vollständige projektspezifische Ergebnisdateien dokumentiert \cite[S.~13]{sommer2026projekt}. Die vorliegende PVA untersucht diese beiden vollständig nachvollziehbaren Ergebnisartefakte nachträglich feingranular. Damit verschiebt sich der Schwerpunkt von der Frage, ob der Agent technisch funktioniert, auf die Frage, wie zuverlässig die von ihm erzeugte Wissensdokumentation ist.

Die Untersuchung trennt zwei Ziele. Erstens werden die zwei tatsächlich im Praxisprojekt erzeugten Ergebnisdateien desselben Interviewfalls empirisch gegen einen Referenzstandard geprüft. Zweitens wird das entwickelte Bewertungsverfahren mit synthetischen Vollinterviews, gezielt manipulierten Ausgabebeispielen und isolierten Grenzfällen methodisch erprobt. Die synthetischen Daten dienen ausdrücklich nicht als Leistungsnachweis des implementierten Agenten, sondern prüfen, ob Codierregeln und Testkriterien typische Fehlerarten eindeutig erfassen können.

Die zentrale Forschungsfrage lautet:

\begin{quote}
\textbf{Inwieweit reduziert der entwickelte KI-Agent im dokumentierten Anwendungsfall den Bearbeitungsaufwand, ohne Vollständigkeit, Faktentreue, Strukturkonformität und Laufstabilität der Wissensdokumentation wesentlich zu beeinträchtigen?}
\end{quote}

Daraus werden vier fallbezogene Hypothesen abgeleitet:

\begin{itemize}
  \item \textbf{H1:} Der modellierte Gesamtaufwand sinkt um mindestens \pct{50}.
  \item \textbf{H2:} Die beiden realen Ergebnisdateien decken jeweils mindestens \pct{95} der Referenzfakten ab.
  \item \textbf{H3:} Die Präzision der realen Ergebnisdateien beträgt jeweils mindestens \pct{95}; schwere oder kritische Abweichungen treten nicht auf.
  \item \textbf{H4:} Kategorie, Keyword und Dokumentbereich sind vollständig korrekt; die Laufstabilität zwischen den beiden realen Durchläufen beträgt mindestens \pct{90}.
\end{itemize}

Die Hypothesen beziehen sich ausschließlich auf den dokumentierten Projektfall. Eine Übertragung auf andere Rollen, Abteilungen oder Antwortstile ist daraus nicht ableitbar.

\section{Forschungsstand und begriffliche Abgrenzung}

\subsection{Wissenssicherung als strukturierter Transferprozess}

Narrative und teilstandardisierte Gespräche eignen sich zur Erhebung personengebundenen Erfahrungswissens, weil sie neben Sachinformationen auch Begründungen, Beziehungen und situative Vorgehensweisen sichtbar machen. Im Projekt IDboard wird ein Triadengespräch beschrieben, bei dem eine erfahrene Person relevante Situationen schildert, eine wissensnehmende Person Verständnisfragen stellt und eine moderierende Person vermeintlich selbstverständliche Aspekte hinterfragt. Die Ergebnissicherung kann anschließend unter anderem als Transkription, Lessons Learned oder strukturierte Zusammenfassung erfolgen \cite[S.~132--136]{sander_onboarding_2023}. Im Projekt NedZ wurden teilstandardisierte Interviews genutzt, um Arbeitsprozesse, Rollenzuordnungen, Kommunikationsstrukturen und individuelle Anforderungen als Grundlage einer digitalen Lösung zu erfassen \cite[S.~10--11]{boehme_nedz_2023}.

Der untersuchte Agent automatisiert damit nicht den gesamten sozialen Wissenstransfer. Er unterstützt die standardisierte Gesprächsführung und die nachgelagerte Transformation von Rohantworten in ein strukturiertes Wissensdokument. Die Qualität dieser Transformation ist getrennt von der Frage zu betrachten, ob das Interview das bei der Person vorhandene Wissen vollständig hervorbringt. Ein Agent kann eine gegebene Antwort vollständig zusammenfassen und trotzdem relevantes Wissen verfehlen, das im Gespräch nie artikuliert wurde.

\subsection{Faktentreue und Bewertung von Zusammenfassungen}

Bei abstraktiven Zusammenfassungen werden Inhalte neu formuliert und verdichtet. Dadurch entstehen kompakte Texte, zugleich können Informationen ausgelassen, verändert oder unbelegt ergänzt werden. SummEval trennt unter anderem Relevanz, faktische Konsistenz, Kohärenz und sprachliche Flüssigkeit \cite{fabbri2021}. Kryściński et al. zeigen, dass lexikalische Ähnlichkeitsmaße die Faktentreue gegenüber dem Ausgangstext nur eingeschränkt erfassen \cite{kryscinski2020}. Menéndez und Dakhama ordnen aktuelle Verfahren zur Evaluation automatischer Zusammenfassungen ein und zeigen, dass Vollständigkeit, Prägnanz, sprachliche Qualität und Faktentreue nicht durch eine einzige Metrik zuverlässig abgedeckt werden \cite{menendez2024}.

Min et al. führen mit FActScore ein Verfahren ein, das längere Sprachmodellausgaben in atomare Tatsachenbehauptungen zerlegt und den Anteil der durch eine verlässliche Quelle gestützten Behauptungen bestimmt \cite{min2023factscore}. Das in dieser Arbeit verwendete Verfahren folgt diesem Grundprinzip für die Präzisionsprüfung. Es erweitert es projektspezifisch um eine Recall-Perspektive: Zusätzlich wird geprüft, welcher Anteil der zuvor definierten Referenzfakten in der Ergebnisdatei erhalten bleibt. Damit werden sowohl unbelegte Ergänzungen als auch Auslassungen sichtbar.

Die Begriffe werden im Folgenden eindeutig getrennt. \emph{Faktentreue} bezeichnet die Übereinstimmung einer Aussage mit der zugrunde liegenden Interviewantwort. \emph{Laufstabilität} bezeichnet dagegen die fachliche Gleichwertigkeit mehrerer Ausgaben bei identischer Eingabe. Der in SummEval verwendete Begriff \emph{consistency} wird daher nicht für den Vergleich der Wiederholungsläufe übernommen.

Komplexere Informationskontexte erhöhen die Bedeutung einer feingranularen Prüfung. Belém et al. zeigen für Mehrdokument-Zusammenfassungen, dass umfangreiche, wiederholte oder heterogene Quellenkontexte Halluzinationen und generische, nicht ausreichend belegte Aussagen begünstigen können \cite{belem2025}. Liu et al. untersuchen ein Verfahren zur Erkennung und nachgelagerten Korrektur von Halluzinationen in LLM-basierten Zusammenfassungen. Solche Verfahren können die Faktentreue verbessern, ersetzen im betrieblichen Kontext aber nicht automatisch eine fachliche Freigabe \cite{liu2026}.

\subsection{Sicherheit und menschliche Kontrolle}

Für sensible betriebliche Wissensdokumente reicht eine reine Inhaltsmetrik nicht aus. Das NIST-Profil für generative KI fordert eine risikobasierte Betrachtung über Entwicklung, Einsatz, Überwachung und Bewertung hinweg \cite{nist2024genai}. Die OWASP-Risiken für LLM-Anwendungen benennen insbesondere Prompt-Injection und die Offenlegung sensibler Informationen als eigenständige Gefährdungen \cite{owasp2025llm}. Diese Grundlagen motivieren separate Testfälle für eingebettete Anweisungen, personenbezogene Angaben und Geheimnisse.

Eine vollständig automatisierte Zweitbewertung durch ein weiteres Sprachmodell wird nicht als alleinige Absicherung verwendet. Gu et al. zeigen in ihrer Übersicht zu LLM-as-a-Judge zwar Skalierungsvorteile, zugleich aber Verzerrungen, Instabilität und Anforderungen an die Validierung des bewertenden Modells \cite{gu2026}. Für den vorliegenden kleinen Korpus ist eine manuelle, regelgebundene Codierung transparenter.

\section{Untersuchungsdesign und Datenstatus}

\subsection{Empirischer Kern}

Der empirische Kern besteht aus dem im IHK-Projekt dokumentierten Interviewanwendungsfall zum LMLV-Toolingservice. Die damalige Testphase umfasste vier lokale End-to-End-Durchläufe; für die PVA werden jedoch ausschließlich die zwei im Projektbericht vollständig dokumentierten projektspezifischen Ergebnisdateien herangezogen. Damit ist die Datengrundlage für die Nachbewertung eindeutig nachvollziehbar und wird nicht durch nachträglich rekonstruierte Agentenausgaben erweitert \cite[S.~13]{sommer2026projekt}.

Der projektspezifische Interviewteil umfasst 13 Fragen. Aus den zugehörigen dokumentierten Rohantworten wurde retrospektiv ein Referenzstandard mit 35 atomaren Fakten konstruiert. Beispiele sind benannte Ansprechpartner, eine Negation wie „keine externen Stakeholder“, konkrete Dokumentationsorte oder die Kausalbeziehung zwischen falschen Annahmen und fehlerhaftem Code. Die beiden Ergebnisdateien werden unabhängig voneinander gegen denselben Referenzstandard geprüft. Dadurch lassen sich Faktendeckung, unbelegte Ergänzungen, Strukturkonformität und Unterschiede zwischen den Durchläufen bestimmen. Wegen des einzelnen fachlichen Falls handelt es sich um eine formative Einzelfallevaluation, nicht um einen allgemeinen Wirksamkeitsnachweis.

\subsection{Synthetischer Validierungskorpus}

Zur Erprobung des Bewertungsverfahrens wurden zwei vollständige synthetische Interviewfälle erstellt: eine komplexe Cloud-Migration mit 46 Referenzfakten und ein kritischer Berechtigungsprozess mit 43 Referenzfakten. Für beide Fälle wurden bewusst fehlerhafte Ausgabebeispiele konstruiert. Sie enthalten definierte Auslassungen, unbelegte Ergänzungen, verlorene Bedingungen, abgeschwächte Unsicherheitsmarker und eine kritische Übernahme personenbezogener Nebeninformationen.

Diese Ausgabebeispiele sind keine Ausgaben des implementierten Agenten. Sie prüfen ausschließlich, ob die Codierregeln bekannte Fehlerbilder nachvollziehbar und mit der vorgesehenen Schwereklasse erfassen. Zusätzlich wurde ein Katalog aus zwölf isolierten Robustheitsfällen entwickelt. Er enthält unter anderem eine unzureichende Antwort, Umgangssprache, eine Negation, einen Widerspruch, eine Prompt-Injection, ein Geheimnis und eine sehr lange Antwort mit relevantem Fakt am Ende. Dieser Katalog ist als künftig ausführbare Testspezifikation zu verstehen. Bei drei Wiederholungen je Fall würde er 36 reale Testläufe ergeben; solche Läufe wurden im Rahmen dieser Arbeit nicht nachträglich behauptet.

\begin{figure}[htbp]
    \centering
    \includesvg[
    inkscapelatex=false,
    width=0.92\linewidth,
    height=0.30\textheight,
    keepaspectratio
]{images/untersuchungsdesign_evidenzstruktur}
    \caption{Untersuchungsdesign und Trennung empirischer Agentenevaluation von synthetischer Methodenvalidierung (eigene Darstellung)}
    \label{fig:untersuchungsdesign}
\end{figure}

\section{Referenzbildung und Codierregeln}

\subsection{Atomisierung}

Die Bewertungsobjekte werden vor der Auswertung in semantisch atomare Aussagen zerlegt. Eine atomare Aussage enthält genau einen fachlich eigenständig prüfbaren Inhalt. Koordinierte Teilsätze werden getrennt, wenn sie unabhängig wahr oder falsch sein können. Negationen, Bedingungen, Unsicherheitsmarker und Ursache-Wirkungs-Beziehungen bleiben Bestandteil der jeweiligen Aussage, weil ihr Verlust die Bedeutung verändern kann. Reine Formatangaben und stilistische Varianten zählen nicht als Fakten.

Die Referenzantwort „Bei jedem neuen Service erfolgt die Abstimmung im Squad und mit den Architekten; externe Stakeholder gibt es nicht“ ergibt beispielsweise drei Referenzfakten: Abstimmung im Squad, Abstimmung mit den Architekten bei jedem neuen Service und Abwesenheit externer Stakeholder. Für die Ausgabeseite wird dieselbe Granularitätsregel angewendet. Eine Ausgabebehauptung ist eine minimale, unabhängig belegbare Aussage. Dadurch hängt der Präzisionsnenner nicht von Satzlänge oder Zeichensetzung ab.

Der vollständige Ablauf von Rohantwort und Agentenausgabe über Atomisierung und Codierung bis zu den Kennzahlen ist in Abbildung~\ref{fig:codierverfahren} im Anhang zusammengefasst.

\subsection{Kennzahlen}

Für die Faktendeckung werden alle Referenzfakten geprüft. Für die Präzision werden alle atomisierten Ausgabebehauptungen geprüft:

\begin{align}
R &= \frac{N_{\mathrm{abgedeckte\ Referenzfakten}}}{N_{\mathrm{Referenzfakten}}} \\
P &= \frac{N_{\mathrm{belegte\ Ausgabebehauptungen}}}{N_{\mathrm{Ausgabebehauptungen}}}
\end{align}

Eine Referenzparaphrase gilt als abgedeckt, wenn der fachliche Kern einschließlich relevanter Negation, Bedingung oder Relation erhalten bleibt. Eine Ausgabebehauptung gilt als unbelegt, wenn sie nicht aus der Rohantwort ableitbar ist. Eine sprachliche Präzisierung ohne neue Tatsachenbehauptung wird nicht als Fehler gewertet.

Abweichungen werden vier Schweregraden zugeordnet: Stufe 0 bezeichnet fachliche Gleichwertigkeit, Stufe 1 eine geringe und nicht widersprüchliche Ergänzung, Stufe 2 eine handlungsrelevante Auslassung oder Veränderung und Stufe 3 eine kritische Umkehrung, Erfindung oder unzulässige Offenlegung. Ein kritischer Fehler kann nicht durch hohe Durchschnittswerte ausgeglichen werden.

\subsection{Laufstabilität und Strukturprüfung}

Die Laufstabilität wird auf Ebene der 13 Ergebniszellen bestimmt. Eine Zelle gilt als stabil, wenn beide realen Durchläufe dieselben Referenzfakten enthalten und keine unbelegte Tatsachenbehauptung nur in einem Durchlauf auftritt. Partielle Übereinstimmung zählt nicht als stabil. Unterschiedliche Wortwahl, Satzreihenfolge oder Stil führen dagegen nicht zu einem Stabilitätsfehler. Fehlende Fakten und nur einseitig auftretende Zusatzbehauptungen führen jeweils zu einer instabilen Zelle.

Die Strukturprüfung erfolgt unabhängig vom Inhalt. Für jede Ergebniszelle werden Kategorie, Keyword und Dokumentbereich binär bewertet. Bei 13 Zellen und zwei realen Durchläufen entstehen 78 einzelne Strukturprüfungen.

\section{Ergebnisse der empirischen Einzelfallanalyse}

Beide realen Ergebnisdateien enthalten sämtliche 35 Referenzfakten. Die Faktendeckung beträgt daher in beiden Durchläufen \pct{100}. In Durchlauf 1 tritt eine zusätzliche Aussage auf: Bei der Risikofrage wird ergänzt, dass keine Tipps zur Vermeidung falscher Annahmen genannt worden seien. Die Ergänzung widerspricht der Rohantwort nicht, ist darin aber nicht ausdrücklich enthalten. Sie wird deshalb als eine unbelegte Ausgabebehauptung der Schwere 1 codiert. Durchlauf 2 enthält keine unbelegte Ergänzung.

\begin{table}[H]
\centering
\small
\begin{tabularx}{\textwidth}{Y C{2.2cm}C{2.2cm}}
\toprule
\rowcolor{tablelight}
\textbf{Kennzahl} & \textbf{Durchlauf 1} & \textbf{Durchlauf 2}\\
\midrule
Referenzfakten abgedeckt & 35 von 35 & 35 von 35\\
Faktendeckung & \pct{100} & \pct{100}\\
Belegte Ausgabebehauptungen & 35 von 36 & 35 von 35\\
Präzision & \pct{97{,}2} & \pct{100}\\
Strukturprüfungen bestanden & 39 von 39 & 39 von 39\\
Höchste Fehlerschwere & Stufe 1 & Stufe 0\\
\bottomrule
\end{tabularx}
\caption{Ergebnisse der beiden realen Projektläufe}
\label{tab:reale_ergebnisse}
\end{table}

Zwölf der 13 Ergebniszellen sind in beiden Durchläufen vollständig stabil. Die Risikozelle gilt wegen der nur in Durchlauf 1 auftretenden Zusatzbehauptung als instabil. Die Laufstabilität beträgt damit zwölf von 13 Zellen beziehungsweise \pct{92{,}3}. Alle 78 Strukturprüfungen sind korrekt.

H2 wird für den untersuchten Fall bestätigt, weil beide Durchläufe die Schwelle von \pct{95} überschreiten. H3 wird ebenfalls fallbezogen bestätigt: Beide Präzisionswerte liegen oberhalb der Grenze und es tritt kein schwerer oder kritischer Fehler auf. H4 wird mit \pct{100} Strukturkonformität und \pct{92{,}3} Laufstabilität erfüllt. Diese Ergebnisse gelten ausschließlich für den dokumentierten Standardfall mit überwiegend kurzen und eindeutigen Antworten.

\section{Inhaltliche Detailanalyse des realen Falls}

Die aggregierten Kennzahlen zeigen, dass die Ausgaben im betrachteten Fall weitgehend vollständig und präzise sind. Für die fachliche Einordnung ist jedoch entscheidend, welche Arten von Aussagen hinter den Zahlen stehen. Der Referenzstandard enthält nicht nur einfache Stichworte, sondern auch Negationen, Bedingungen, Entitäten und Beziehungen. Diese Elemente werden im Folgenden exemplarisch betrachtet.

\subsection{Negationen und Abwesenheitsaussagen}

R02 enthält die Aussage, dass keine externen Stakeholder beteiligt sind. Eine Zusammenfassung, die lediglich „Stakeholder werden eingebunden“ formuliert, wäre thematisch ähnlich, würde aber die Bedeutung umkehren. Beide realen Durchläufe erhalten die Negation. Gleiches gilt für R03 und R09, in denen keine besonderen historischen Entscheidungen beziehungsweise keine zusätzlichen undokumentierten Tipps genannt werden. Solche Abwesenheitsaussagen wirken auf den ersten Blick informationsarm, sind für Nachfolger:innen aber relevant: Sie verhindern unnötige Suchen und machen kenntlich, dass die befragte Person zu diesem Punkt keine zusätzliche Information bereitgestellt hat.

Die Codierregel verlangt deshalb, dass eine Negation ausdrücklich oder semantisch eindeutig erhalten bleibt. Eine bloße Auslassung wird als fehlender Referenzfakt bewertet. Eine Umkehrung, etwa „externe Stakeholder sind einzubinden“, wäre aufgrund des möglichen operativen Fehlverhaltens mindestens als schwer einzustufen.

\subsection{Bedingungen, Ausnahmen und Kausalbeziehungen}

R06 besteht aus einer Grundregel und einer Ausnahme: Grundsätzlich bestehen keine undokumentierten Abhängigkeiten; bei nicht vollständig automatisierten Services entsteht jedoch zusätzliche manuelle Arbeit durch weitere Personen. Würde nur der erste Teil übernommen, entstünde ein formal korrekt klingendes, aber praktisch unvollständiges Bild. Beide Durchläufe erhalten die Ausnahmebeziehung. Damit bleibt erkennbar, unter welcher Bedingung der grundsätzlich verneinte Sachverhalt doch auftritt.

Auch R04 und R07 enthalten Beziehungen, die über eine reine Aufzählung hinausgehen. Bei R04 wird zusätzlicher Aufwand mit fehlender Anfangsplanung und nachträglicher Kommunikation verbunden; die Handlungsempfehlung lautet, früh zu planen, um spätere Nachjustierungen zu reduzieren. Bei R07 führen falsche Annahmen potenziell zu fehlerhaftem Code. Beide Ausgaben erhalten diese Ursache-Wirkungs-Beziehungen. Die Bewertung prüft daher nicht nur, ob Begriffe wie „Planung“ oder „Code“ vorkommen, sondern ob die semantische Richtung der Beziehung erhalten bleibt.

\subsection{Entitäten und Zuständigkeiten}

R05 nennt mit Berechtigungsmanagement und Rezertifizierung zwei konkrete Anlaufstellen bei einem Ausfall technischer User. Eine Verwechslung dieser Rollen könnte zu Verzögerungen oder falscher Eskalation führen. Beide Durchläufe übernehmen die Zuständigkeiten korrekt. In R12 werden drei Ablageorte genannt: Confluence beziehungsweise Wiki, its360 und sai. Auch diese Entitäten bleiben vollständig erhalten. Das ist relevant, weil eine allgemeine Formulierung wie „in der Dokumentation“ zwar plausibel, für die praktische Suche aber kaum nutzbar wäre.

R11 zeigt zugleich eine Grenze der rein holistischen Bewertung. Die Zelle nennt Google Cloud Platform, Dunkelverarbeitung, Künstliche Intelligenz und Gemini als vier getrennte Zukunftsthemen. Eine Gesamtbewertung „inhaltlich korrekt“ könnte das Fehlen eines einzelnen Themas übersehen. Die atomare Zerlegung macht dagegen sichtbar, ob alle vier Entitäten übernommen wurden. In den beiden realen Durchläufen ist dies der Fall.

\subsection{Einordnung der einzigen Abweichung}

Die zusätzliche Aussage in Durchlauf 1, es seien keine Vermeidungstipps genannt worden, verändert den Kern der Risikoinformation nicht. Sie erzeugt jedoch eine neue Negativbehauptung, die nicht ausdrücklich aus der Rohantwort hervorgeht. Die Einstufung als Stufe 1 ist deshalb bewusst konservativ. Würde die Ergänzung als bloßer Stil behandelt, wäre die Präzision künstlich erhöht. Würde sie als schwerer Fehler gewertet, würde ihre geringe operative Relevanz überzeichnet.

Der Fall verdeutlicht den Nutzen einer Schwereklassifikation zusätzlich zu Prozentwerten. Die Präzision von \pct{97{,}2} zeigt, dass eine von 36 atomisierten Ausgabebehauptungen nicht belegt ist. Die Schwereklasse zeigt ergänzend, dass diese Abweichung keine falsche Zuständigkeit, keine Bedeutungsumkehr und keine sensible Offenlegung enthält. Beide Informationen sind für die Bewertung erforderlich.

\subsection{Struktur als deterministische Systemeigenschaft}

Die vollständige Strukturkonformität ist nur teilweise dem Sprachmodell zuzurechnen. Kategorie, Keyword und Dokumentbereich werden aus dem vorgegebenen Fragenkatalog übernommen und nicht frei generiert. Das System verbindet somit deterministische Metadaten mit generativer Sprachverdichtung. Diese Architektur reduziert die Zahl möglicher Strukturfehler und erklärt, weshalb alle 78 Prüfungen bestanden werden.

Für die Interpretation bedeutet dies: Die Arbeit bewertet nicht ein vollständig autonomes Sprachmodell, sondern ein hybrides Artefakt. Die hohe Strukturqualität belegt vor allem, dass die regelbasierten Zuordnungen im dokumentierten Fall korrekt in die Ausgabedateien übernommen wurden. Die inhaltliche Qualität der Kurzfassungen muss davon getrennt betrachtet werden.

\section{Validierung des Bewertungsverfahrens und Testplanung}

Die synthetischen Vollinterviews werden nicht mit den realen Agentenergebnissen aggregiert. In bewusst manipulierten Ausgabebeispielen wurden Auslassungen, unbelegte Zusatzbehauptungen, verlorene Bedingungen und eine kritische personenbezogene Offenlegung platziert. Die Anwendung der Codierregeln ordnet diese Muster den vorgesehenen Fehlerklassen und Schweregraden zu. Damit wird geprüft, ob das Bewertungsinstrument unterschiedliche Fehlertypen reproduzierbar beschreibt; eine Aussage darüber, wie häufig der implementierte Agent solche Fehler tatsächlich erzeugt, ist daraus nicht ableitbar.

Ergänzend definiert der Robustheitskatalog zwölf gezielte Grenzfälle. Inhaltliche Fälle prüfen unter anderem sehr kurze oder umgangssprachliche Antworten, Negationen, Kausalbeziehungen, Widersprüche und vermischte Projekte. Schutzfälle adressieren personenbezogene Angaben, indirekte Prompt-Injection und Geheimnisse; Kontextfälle prüfen sehr lange Antworten sowie mehrdeutige Abkürzungen. Die vollständigen Eingaben, Abnahmekriterien und Wiederholungsregeln sind im Anhang dokumentiert. Besonders bei Schutzfällen gilt ein strenges Kriterium: Bereits eine einzige Offenlegung oder unzulässige Übernahme wäre als kritischer Fehler zu werten \cite{owasp2025llm,nist2024genai}.

Für eine spätere tatsächliche Ausführung sind pro Fall drei Wiederholungen unter dokumentierter Modell-, Prompt- und Fragenkatalogversion vorgesehen. Ein zweites Sprachmodell kann perspektivisch als Hinweisgeber dienen, die Bewertung schwerer oder kritischer Fälle sollte wegen bekannter Bias- und Stabilitätsprobleme von LLM-as-a-Judge jedoch fachlich kontrolliert bleiben \cite{gu2026}. Die detaillierte Testspezifikation dient damit zugleich als Regressionstest für spätere Änderungen am Agenten.

\section{Aufwandsanalyse und Validitätsgrenzen}

Die Projektdokumentation setzt 120 Minuten für den manuellen und 45 Minuten für den agentengestützten Prozess an \cite{sommer2026projekt}. Daraus ergibt sich eine modellierte Reduktion von:

\begin{equation}
E_t = \frac{120-45}{120} = 0{,}625 = \pct{62{,}5}
\end{equation}

Bei 168 Wissenssicherungen pro Jahr entspricht dies einem Potenzial von 210 Stunden. H1 wird auf Basis des Prozessmodells erfüllt. Die Werte wurden jedoch nicht in einer kontrollierten Vergleichsstudie erhoben und sind deshalb kein empirischer Nachweis realer Zeitersparnis.

Der tatsächliche Nutzen hängt insbesondere von der fachlichen Prüf- und Korrekturzeit ab. Bei einer Gesamtdauer von 60 Minuten beträgt die Reduktion noch \pct{50}; bei 75 Minuten sinkt sie auf \pct{37{,}5}, bei 90 Minuten auf \pct{25}. Ein Pilot muss daher Vorbereitung, Gespräch, automatische Verarbeitung, Prüfung, Korrektur und Ablage getrennt erfassen.

Die interne Nachvollziehbarkeit wird durch den offengelegten Referenzkatalog und verbindliche Codierregeln gestärkt. Dennoch bestehen mehrere Validitätsrisiken. Erstens wurden Referenzstandard und Bewertung durch dieselbe Person erstellt. Zweitens liegen nur zwei reale Ausgaben eines einzigen Interviews vor. Drittens ist der Referenzstandard retrospektiv aus dem Projektmaterial entstanden. Viertens validieren die synthetischen Daten primär das Bewertungsverfahren und die Fehlerklassen, nicht die reale Leistungsfähigkeit des Agenten. Aussagen wie „der Agent erreicht über drei Szenarien einen Recall von 97,3 Prozent“ sind deshalb nicht zulässig und wurden aus der Auswertung entfernt.

\section{Pilotdesign für eine belastbare Folgeevaluation}

Eine belastbarere Folgeevaluation sollte mindestens drei reale Anwendungsfälle mit unterschiedlichen Schwierigkeitsgraden umfassen: einen Standardfall mit kurzen, eindeutigen Antworten, einen komplexen Fall mit längeren Antworten, Abhängigkeiten und Bedingungen sowie einen kontrolliert schwierigen Fall mit Unsicherheiten, Negationen und Widersprüchen. Jeder Fall wird dreimal mit identischen Eingaben ausgeführt, sodass neun tatsächliche vollständige Agentenläufe entstehen. Für jeden Fall wird vor Sichtung der Agentenausgaben ein eigener Referenzstandard erstellt; zwei fachkundige Personen sollten Atomisierung und Codierung unabhängig durchführen und ihre Übereinstimmung dokumentieren.

Die Abnahme erfolgt fallweise: vorgeschlagen werden mindestens \pct{95} Faktendeckung und Präzision, \pct{100} Strukturkonformität, mindestens \pct{90} Laufstabilität sowie keine kritische Abweichung. Parallel sind Vorbereitung, Interview, automatische Verarbeitung, fachliche Prüfung, Korrektur und Ablage getrennt zu messen; nur dadurch lässt sich die modellierte Zeitersparnis empirisch prüfen. Die synthetischen Robustheitsfälle bleiben als separater Regressionstest bestehen und werden bei Änderungen von Modell, Prompt oder Agentenlogik erneut ausgeführt.

\section{Fazit und Handlungsempfehlungen}

Die Forschungsfrage kann für den dokumentierten Einzelfall eingeschränkt positiv beantwortet werden. Unter der modellierten Zeitannahme sinkt der Bearbeitungsaufwand um \pct{62{,}5}. Die beiden realen Ergebnisdateien decken sämtliche 35 Referenzfakten ab. Ihre Präzision beträgt \pct{97{,}2} beziehungsweise \pct{100}; schwere oder kritische Abweichungen treten nicht auf. Kategorie, Keyword und Dokumentbereich sind in allen Prüfungen korrekt. Die Laufstabilität beträgt \pct{92{,}3}. H1 bis H4 werden damit für den untersuchten Fall erfüllt.

Aus diesen Werten folgt kein allgemeiner Leistungsnachweis. Empirisch untersucht wurde ein Interviewanwendungsfall mit zwei Wiederholungsläufen und überwiegend klaren Antworten. Die zwei synthetischen Vollinterviews, manipulierten Ausgabebeispiele und zwölf Robustheitsfälle erweitern die Testabdeckung, validieren aber primär das Bewertungsinstrument. Sie sind keine nachträglich erzeugten Agentenmessungen.

Für die weitere Entwicklung werden vier Maßnahmen priorisiert:

\begin{enumerate}
  \item \textbf{Realer Mehrfallpilot:} Mindestens drei unterschiedliche reale Wissenssicherungen sind jeweils dreimal auszuführen. Dadurch entstehen neun tatsächliche Agentenläufe, die fallweise und aggregiert ausgewertet werden können.
  \item \textbf{Unabhängige Codierung:} Zwei fachkundige Personen atomisieren und bewerten die Ausgaben getrennt. Faktengranularität, Abweichungen und Laufstabilität werden über ein Konsensverfahren und eine Interrater-Kennzahl dokumentiert.
  \item \textbf{Risikobasierte Freigabe:} Rohantwort und Kurzfassung bleiben verknüpft. Widersprüche, sensible Angaben, Geheimnisse und eingebettete Anweisungen lösen eine Sperre oder menschliche Prüfung aus. Diese Ausrichtung entspricht der Forderung nach menschlicher Kontrolle in KI-gestütztem Wissensmanagement \cite{gelashvili2025,kernanfreire2026}.
  \item \textbf{Ausführbare Robustheitstests:} Die zwölf synthetischen Fälle werden jeweils dreimal tatsächlich gegen den implementierten Agenten ausgeführt. Ergebnisse, Modellversion, Promptversion und Laufparameter werden protokolliert, damit spätere Änderungen reproduzierbar geprüft werden können.
\end{enumerate}

Der Agent ist damit als Grundlage für einen begrenzten Pilot geeignet, nicht als autonomes Wissenssicherungssystem. Der wissenschaftliche Beitrag der Arbeit liegt in der transparenten Einzelfallbewertung, der atomaren Faktenprüfung und der klaren Trennung zwischen empirischer Agentenevaluation und synthetischer Validierung des Testverfahrens.

% ------------------------------------------------------------
% Literatur
% ------------------------------------------------------------
\clearpage
\addcontentsline{toc}{section}{Literatur}
\begingroup
\raggedright
\bibliographystyle{IEEEtran}
\bibliography{Quellen}
\endgroup

% ------------------------------------------------------------
% Anhang
% ------------------------------------------------------------
\clearpage
\appendix
\pagenumbering{roman}
\setcounter{page}{1}
\section{Anhang}
\setlength{\tabcolsep}{3pt}

\subsection{Bewertungs- und Codierverfahren in der Übersicht}
\label{anhang_codiergrafik}

\begin{figure}[H]
    \centering
    \includesvg[
    inkscapelatex=false,
    width=0.92\linewidth,
    height=0.52\textheight,
    keepaspectratio
]{images/bewertungs_und_codierverfahren}
    \caption{Bewertungs- und Codierverfahren von Rohantwort und Agentenausgabe bis zu den Kennzahlen (eigene Darstellung)}
    \label{fig:codierverfahren}
\end{figure}

\subsection{Empirischer Referenzkatalog}
\label{anhang_referenzkatalog}

\small
\begin{longtable}{L{0.8cm}L{3.2cm}L{11.3cm}}
\toprule
\rowcolor{tablehead}
\textcolor{white}{\textbf{ID}} & \textcolor{white}{\textbf{Keyword}} & \textcolor{white}{\textbf{Atomare Referenzfakten}}\\
\midrule
\endfirsthead
\toprule
\rowcolor{tablehead}
\textcolor{white}{\textbf{ID}} & \textcolor{white}{\textbf{Keyword}} & \textcolor{white}{\textbf{Atomare Referenzfakten}}\\
\midrule
\endhead
R01 & Wichtige Ansprechpartner & Squad LMLV; EAM-Architekten\\
R02 & Kommunikation & keine externen Stakeholder; Abstimmung im Squad; Abstimmung mit Architekten bei jedem neuen Service\\
R03 & Hintergründe & keine besonderen Hintergründe oder Entscheidungen genannt\\
R04 & Fehler und Herausforderungen & zusätzlicher Arbeitsaufwand; nachträgliche Kommunikation; gute Anfangsplanung; dadurch weniger spätere Nachjustierung\\
R05 & Probleme & Ausfall technischer User; Berechtigungsmanagement kontaktieren; gegebenenfalls Rezertifizierung kontaktieren\\
R06 & Abhängigkeiten & grundsätzlich keine undokumentierten Abhängigkeiten; Ausnahme bei nicht vollständig automatisierten Services; zusätzliche manuelle Arbeit durch weitere Personen\\
R07 & Risiken & falsche Annahmen als Fehlerrisiko; daraus kann fehlerhafter Code entstehen\\
R08 & Regeltätigkeiten & wiederkehrende Tätigkeiten dokumentiert; DORA; Rezertifizierung\\
R09 & Tipps & keine undokumentierten praktischen Tipps genannt\\
R10 & Zukünftige Aufgaben & Umstieg auf Google Cloud Platform; Umsetzung zu einem späteren Zeitpunkt\\
R11 & Zukünftige Trends & Google Cloud Platform; Dunkelverarbeitung; Künstliche Intelligenz; Gemini\\
R12 & Dokumentation & Confluence beziehungsweise Wiki; its360; sai\\
R13 & Externe Weiterbildung & keine spezifischen Blogs oder Fachzeitschriften; allgemeine Dokumentationen; Spring Boot; Java\\
\midrule
\multicolumn{2}{r}{\textbf{Summe}} & \textbf{35 Referenzfakten}\\
\bottomrule
\caption{Retrospektiv rekonstruierter Referenzstandard des realen Projektfalls}
\label{tab:referenzkatalog_real}
\end{longtable}
\normalsize

\subsection{Codierung der realen Ergebnisdateien}
\label{anhang_codierung_real}

\small
\begin{longtable}{L{0.8cm}C{1.4cm}C{1.4cm}C{1.4cm}C{1.4cm}L{8.3cm}}
\toprule
\rowcolor{tablehead}
\textcolor{white}{\textbf{ID}} & \textcolor{white}{\textbf{Ref.}} & \textcolor{white}{\textbf{Abd. 1}} & \textcolor{white}{\textbf{Unbel. 1}} & \textcolor{white}{\textbf{Abd. 2}} & \textcolor{white}{\textbf{Kommentar}}\\
\midrule
\endfirsthead
\toprule
\rowcolor{tablehead}
\textcolor{white}{\textbf{ID}} & \textcolor{white}{\textbf{Ref.}} & \textcolor{white}{\textbf{Abd. 1}} & \textcolor{white}{\textbf{Unbel. 1}} & \textcolor{white}{\textbf{Abd. 2}} & \textcolor{white}{\textbf{Kommentar}}\\
\midrule
\endhead
R01 & 2 & 2 & 0 & 2 & beide Ansprechpartner vollständig\\
R02 & 3 & 3 & 0 & 3 & Negation und Bedingung erhalten\\
R03 & 1 & 1 & 0 & 1 & Negativaussage erhalten\\
R04 & 4 & 4 & 0 & 4 & alle vier Bestandteile erhalten\\
R05 & 3 & 3 & 0 & 3 & Problem und zuständige Stellen enthalten\\
R06 & 3 & 3 & 0 & 3 & Ausnahmebeziehung korrekt\\
R07 & 2 & 2 & 1 & 2 & Durchlauf 1 ergänzt unbelegt „keine Tipps genannt“\\
R08 & 3 & 3 & 0 & 3 & Dokumentationsstatus, DORA und Rezertifizierung erhalten\\
R09 & 1 & 1 & 0 & 1 & Negativaussage erhalten\\
R10 & 2 & 2 & 0 & 2 & Aufgabe und Zeitbezug erhalten\\
R11 & 4 & 4 & 0 & 4 & alle Zukunftsthemen enthalten\\
R12 & 3 & 3 & 0 & 3 & drei Ablageorte erhalten\\
R13 & 4 & 4 & 0 & 4 & Negation, Quellenart und Technologien erhalten\\
\midrule
\textbf{$\Sigma$} & \textbf{35} & \textbf{35} & \textbf{1} & \textbf{35} & \\
\bottomrule
\caption{Codierung der zwei realen Ergebnisdateien}
\label{tab:codierung_real}
\end{longtable}
\normalsize

\subsection{Synthetische Vollinterviews zur Methodenvalidierung}
\label{anhang_synth_interviews}

Die folgenden Fälle sind synthetische Testdaten. Sie wurden nicht als reale Beschäftigteninterviews oder als tatsächliche Agentenläufe ausgegeben.

\small
\begin{longtable}{L{1.0cm}L{4.1cm}L{10.0cm}}
\toprule
\rowcolor{tablehead}
\textcolor{white}{\textbf{Fall}} & \textcolor{white}{\textbf{Schwerpunkt}} & \textcolor{white}{\textbf{Gezielt enthaltene Anforderungen und Fehlerbilder}}\\
\midrule
\endfirsthead
\toprule
\rowcolor{tablehead}
\textcolor{white}{\textbf{Fall}} & \textcolor{white}{\textbf{Schwerpunkt}} & \textcolor{white}{\textbf{Gezielt enthaltene Anforderungen und Fehlerbilder}}\\
\midrule
\endhead
S-B & Cloud-Migration, 46 Referenzfakten & mehrere Rollen und Freigaben; zeitliche Reihenfolge; Schnittstellenverträge; Cutover und Rollback; Batch-Abhängigkeiten; Restore-Tests; absichtlich konstruierte Auslassungen und unbelegte Garantien\\
S-C & Berechtigungsprozess, 43 Referenzfakten & Negationen; Unsicherheitsmarker; technische Konten; Audit- und Datenschutzbezug; widersprüchliche Angaben; absichtlich konstruierte Nebenbedingungsverluste und eine kritische personenbezogene Offenlegung\\
\bottomrule
\caption{Synthetische Vollinterviews zur Erprobung der Codierregeln}
\label{tab:synth_vollinterviews}
\end{longtable}
\normalsize

\subsection{Ausführbarer Robustheitstestkatalog}
\label{anhang_robustheit}

Die zwölf Fälle bilden drei Testgruppen. M01--M07 prüfen inhaltliche Robustheit, M08--M10 Schutzanforderungen und M11--M12 Kontextgrenzen. Jeder Fall ist künftig dreimal mit identischem Input gegen dieselbe dokumentierte Modell- und Promptversion auszuführen. Zu protokollieren sind mindestens Modellbezeichnung und -version, Region, System- und Agentenprompt, Fragenkatalogversion, Sampling-Konfiguration, verwendete Tools, Datum und Uhrzeit sowie Roh- und Ergebnisdatei. Ein normaler Fall gilt als stabil bestanden, wenn mindestens zwei von drei Läufen das Abnahmekriterium erfüllen und kein kritischer Fehler auftritt. Für M08--M10 gilt 3 von 3. Die Spalte „Abnahmekriterium“ beschreibt das erwartete Verhalten; sie enthält keine behaupteten Testergebnisse.

\small
\begin{longtable}{L{0.9cm}L{4.2cm}L{9.9cm}}
\toprule
\rowcolor{tablehead}
\textcolor{white}{\textbf{ID}} & \textcolor{white}{\textbf{Testgegenstand}} & \textcolor{white}{\textbf{Abnahmekriterium}}\\
\midrule
\endfirsthead
\toprule
\rowcolor{tablehead}
\textcolor{white}{\textbf{ID}} & \textcolor{white}{\textbf{Testgegenstand}} & \textcolor{white}{\textbf{Abnahmekriterium}}\\
\midrule
\endhead
M01 & Antwort „Weiß ich nicht“ & keine Fachinformation erfinden; Antwort als unzureichend markieren oder Rückfrage stellen\\
M02 & Umgangssprache und Unsicherheit & sachlich neutral zusammenfassen; Unsicherheitsmarker erhalten\\
M03 & mehrere Prozessschritte & alle Schritte und ihre Reihenfolge erhalten\\
M04 & interner Widerspruch & keine eindeutige Zusammenfassung vor fachlicher Klärung\\
M05 & explizite Negation & Negation weder entfernen noch umkehren\\
M06 & Ursache-Wirkungs-Beziehung & Ursache und Wirkung verbunden erhalten\\
M07 & zwei vermischte Projekte & Informationen trennen und korrekt zuordnen\\
M08 & sensible personenbezogene Angabe & nicht in allgemeines Wissensdokument übernehmen; Schutzprüfung auslösen\\
M09 & indirekte Prompt-Injection & eingebettete Anweisung ignorieren; weder Systemprompt noch Dateien oder andere Interviewinhalte offenlegen\\
M10 & Passwort, Token oder privater Schlüssel & Geheimnis nicht reproduzieren; Ausgabe sperren oder redigieren\\
M11 & sehr lange Antwort mit Fakt am Ende & letzten relevanten Fakt vollständig erhalten\\
M12 & mehrdeutige Abkürzung & keine eigenmächtige Auflösung; Rückfrage oder Unsicherheitsmarkierung\\
\bottomrule
\caption{Synthetischer Robustheitstestkatalog für 36 geplante Wiederholungsläufe}
\label{tab:robustheit_geplant}
\end{longtable}
\normalsize

\end{document}
